\begin{table}[t]
    \centering
    \small
    \begin{tabular}{@{}lccccc@{}}
        \toprule
        \textbf{Domain} & \textbf{Items} & \textbf{Seed SD (pts)} & \textbf{Seed range (pts)} & \textbf{Item flip rate} & \textbf{SD / binomial SE} \\
        \midrule
        Code & 164 & \textbf{2.13} & \textbf{4.08} & \textbf{26.7\%} & \textbf{0.61} \\
        Reasoning & 400 & 1.16 & 2.21 & 24.2\% & 0.58 \\
        Factual & 398 & 1.14 & 2.18 & 14.4\% & 0.49 \\
        Creative (judge pts) & 100 & 0.05 & 0.09 & -- & -- \\
        \midrule
        \multicolumn{6}{@{}l}{\emph{Matched $n=164$ items (200 subsamples)}} \\
        Code & 164 & \textbf{2.13} & & & \\
        Reasoning & 164 & 1.85 {\scriptsize [1.61, 2.20]} & & & \\
        Factual & 164 & 1.59 {\scriptsize [1.30, 2.05]} & & & \\
        \bottomrule
    \end{tabular}
    \caption{Seed variability under sampling, averaged over 11 models $\times$ 5 sampled strategies. The flip rate is the share of items whose outcome differs across the 3 seeds. Code is the most seed-sensitive domain on every measure (paired Wilcoxon over 55 cells, $p<10^{-5}$; flip rate of code vs.\ reasoning, $p=0.02$). With matched $n$, code's lead over reasoning shrinks to 0.28 points.}
    \label{tab:seed}
\end{table}
